% ---
% id: meat-core/rem:modes-da1
% kind: rem
% title: "The two source modes on the program"
% summary: >-
%   A declared marginal deletes every feasible state with another marginal, so
%   untrusted-source mode is always available and never gives a larger value.
%   At d_A = 1 the two modes coincide, which lets a protocol with no trusted
%   source be declared there; τ_A = 0 is inadmissible.
% prominence: standard
% status: accepted
% uses: [def:admissible-declaration]
% mentions: [def:roundmodel, def:feasible, fact:meat, prop:trivial-alice-system, transmitter/rem:source-mode, transmitter/ass:src, lem:product-marginal]
% assumes: [A-SRC]
% refs:
%   - "arqand2025meat, Eq. (1)"
% provenance: newly-derived
% lean:
% ---

\begin{remark}\label{rem:modes-da1}
\emph{(a) What each mode asserts} is the transmitter's statement
(\xref{transmitter/rem:source-mode}\,(a): in trusted-source mode the
declaration asserts, through A-SRC
(\xref{transmitter/ass:src}), that Alice's register $A$ never passes
through the adversary's hands, so that $\omega_A=\tau_A$ in every round;
in untrusted-source mode it asserts nothing about $A$). What that does to
this module's program is stated here. In untrusted-source mode
($\tau_A=\nomarg$) the adversary prepares $A\Qsys$ jointly
(Definition~\ref{def:roundmodel}), the \tref{feasible set} carries only
$\tr\omega=1$ (Eq.~\eqref{eq:feas}), and Lemma~\ref{lem:product-marginal} is
not invoked at all. A declared $\tau_A$ \emph{deletes} from the feasible set
every state whose $A$-marginal differs, hence every attack that would have
produced one. That is a different \emph{kind} of input from an entry of
$\CList$, and a more demanding one to justify: an entry of $\CList$
restricts what Eve's channel may do to a register she holds, whereas
$\tau_A$ asserts that a register never reached her at all. Declaring one
without the attestation of Assumption~A-SRC overstates the key in the
unsafe direction. Untrusted-source mode is always available and always
more conservative. Precisely: fix the declaration data
$(\hat O_k,\bowtie_k,c_k)$ and compare the feasible set \eqref{eq:feas}
with $\tau_A$ declared against the one with $\tau_A=\nomarg$. Every
$\omega$ feasible in trusted-source mode is feasible in untrusted-source
mode --- from $\omega_A=\tau_A$ we get $\tr\omega=\tr\tau_A=1$, and taking
the trace of the operator entry
$\tr_{\Qsys}[(\idon{A}\otimes\hat O_k)\,\omega]\bowtie_k c_k\tau_A$
gives the scalar entry
$\tr[(\idon{A}\otimes\hat O_k)\,\omega]\bowtie_k c_k$ (using
$\tr\tau_A=1$, and that $X\succeq Y$ implies $\tr X\ge\tr Y$). Hence
the untrusted feasible set is the larger, its infimum the smaller, and
$\hlow$ computed in untrusted-source mode never exceeds the
trusted-mode value. Dropping the marginal therefore loses rate; it
never gains any, and it can never turn an unsound declaration into a
sound-looking one.

\emph{(b) At $d_A=1$ the two modes are the same declaration}
(Proposition~\ref{prop:trivial-alice-system}). This is the
\emph{embedding}: a protocol with no trusted source may be declared at
$d_A=1$ by folding the system the sender measures into the sectors of
$\Qsys$ and her outcome components into $v$, with the
\tref{party partition} naming which components she holds. The
empty-marginal case is therefore the \emph{safe} direction --- it removes
an assumption rather than adding one.

The equivalence is also the accumulation paper's own reading of its
channel conditional entropy: the sentence immediately following
\cite[Eq.~(1)]{arqand2025meat} records that the definitions of earlier
works are recovered ``by setting the $C$ register to be trivial and
omitting the marginal constraint (equivalently, setting $A$ to be
trivial)''. That sentence is about the source's Eq.~(1), not about its
Theorem~4.2a; the derivation of Proposition~\ref{prop:trivial-alice-system}
is this module's own and does not rest on it. The corresponding check on the
import side is Fact~\ref{fact:meat}, locator note~(vii).

\emph{(c) $\tau_A=0$ is inadmissible}, for the two reasons
\xref{transmitter/rem:source-mode}\,(c) gives; on this module's side
Definition~\ref{def:admissible-declaration}\,(ii) refuses it outright, and were the
normalization dropped, $\omega_A=0$ would give $\tr\omega=0$ and hence
$\omega=0$, so $\Feas=\varnothing$ in violation of
Definition~\ref{def:admissible-declaration}\,(vii) and \eqref{eq:hlow-def} would be an
infimum over the empty set. The absent marginal is $\tau_A=\nomarg$; the
trivial marginal is $d_A=1$ with $\tau_A=[1]$; $\tau_A=0$ is neither.
\end{remark}
